\section{Discussion and Limitations}
\label{sec:discussion}

\subsection{Current Evidence and Key Insights}
The empirical results of the Phase 12 evaluation strongly support the core design philosophy of \method: transferring global topological intelligence into a localized routing agent via offline policy distillation. The selective Predictive Risk Switching (PRS) successfully resolves the "deployment gap" by preserving the low-overhead characteristics of geographic routing in stable environments while dynamically adapting to link breaks and routing holes when they arise. Crucially, the local student policy executes without the need for online global coordinate coordination or active multi-hop message-passing overhead.

\subsection{Remaining Evidence Gaps and Limitations}
While the current results demonstrate clear advantages, several limitations must be addressed to elevate the manuscript to submission quality:
\begin{enumerate}
    \item \textbf{Full-Scale Phase 13 Validation}: Although the sub-modules of the P+ extension (link-loss gate, energy-aware tie-breaking, etc.) are implemented, their statistical verification across the full 14 scenarios and five random seeds remains a primary task.
    \item \textbf{Physical Layer Fidelity}: The evaluations were conducted using a localized Python-based FANET simulator (Lite-GLOBE). While the simulator captures queue dynamics and link expiration rates, it simplifies multi-path fading, shadow fading, and MAC-layer collisions. Cross-validation in network-level simulators like ns-3 is required to demonstrate transferability to real-world hardware.
    \item \textbf{Absence of Traditional Ad Hoc Baselines}: The current comparison set focuses on geographic and machine learning-based baselines. Incorporating traditional reactive (AODV) and proactive (OLSR) protocols will provide a more comprehensive context for network-level overhead claims.
    \item \textbf{Control Overhead Metrics}: The current draft reports input byte sizes as a proxy for overhead. A more rigorous approach should measure the actual control bytes transmitted per delivered packet to provide a clear trade-off comparison against message-passing algorithms like DRAMA.
\end{enumerate}

\subsection{Submission and Positioning Strategy}
For high-impact SCI-level journals (e.g., IEEE Transactions on Mobile Computing or IEEE Internet of Things Journal), the paper must be positioned as a practical, deployable policy distillation method rather than a generic multi-agent RL application. The core narrative should emphasize that global graph knowledge is highly beneficial, but only when restricted to training-time privilege. By demonstrating that the resulting local decision rule is low-overhead, selective, and robust to dynamic topology shifts, the paper directly addresses the practical deployment challenges of UAV swarms.

